\documentclass[aspectratio=169,xcolor=dvipsnames]{beamer}

\AtBeginSection[]
{
    \begin{frame}
        \frametitle{Visão Geral}
        \tableofcontents[currentsection]
    \end{frame}
}

\usefonttheme{professionalfonts} % using non standard fonts for beamer
\usefonttheme{serif} % default family is serif
\usepackage{fontspec}
\setmainfont{XCharter}[
    UprightFont    = *-Roman,
    BoldFont       = *-Bold,
    ItalicFont     = *-Italic,
    BoldItalicFont = *-BoldItalic,
]

\usepackage{hyperref}
\usepackage{graphicx} % Allows including images
\usepackage{booktabs} % Allows the use of \toprule, \midrule and \bottomrule in tables
\input{./slides/SimplePlusTheme/beamerthemeSimplePlus.sty}
\input{./slides/SimplePlusTheme/beamerinnerthemeSimplePlus.sty}
\input{./slides/SimplePlusTheme/beamerfontthemeSimplePlus.sty}
\input{./slides/SimplePlusTheme/beamercolorthemeSimplePlus.sty}

\usepackage{listings}
\setbeamercovered{transparent}

\usepackage{biblatex}
\addbibresource{slides/reference.bib}
\renewcommand*{\bibfont}{\footnotesize}
% Tamanho dos slides de "Leituras Recomendadas": maior que a bibliografia
% completa (\bibfont acima, que rende \tiny via a redefinição de \footnotesize)
% para destacar as leituras curadas.
\newcommand{\leiturasize}{\small}

\usepackage[brazilian ]{babel}
\usepackage{csquotes}
\usepackage{pgfplots}
\pgfplotsset{compat=1.18}

\usepackage{emoji}
\usepackage{amsmath}
\usepackage[xcharter,bigdelims,vvarbb]{newtxmath}
% newtxmath's operators font requests the fontspec family `XCharter(0)' in OT1;
% without these substitutions math digits and operator names silently fall back
% to Computer Modern (LaTeX Font Warning: Font shape `OT1/XCharter(0)/m/n' undefined).
\DeclareFontFamily{OT1}{XCharter(0)}{}
\DeclareFontShape{OT1}{XCharter(0)}{m}{n}{<->ssub * XCharter-TLF/m/n}{}
\DeclareFontShape{OT1}{XCharter(0)}{m}{it}{<->ssub * XCharter-TLF/m/it}{}
\DeclareFontShape{OT1}{XCharter(0)}{b}{n}{<->ssub * XCharter-TLF/b/n}{}
\DeclareFontShape{OT1}{XCharter(0)}{bx}{n}{<->ssub * XCharter-TLF/b/n}{}

\usepackage{bm}
\usepackage[table]{xcolor}
\usepackage{xcolor}
\usepackage{algpseudocode}

\usepackage{cancel}
\usepackage{ulem}

\usetikzlibrary{shapes}
\usetikzlibrary{arrows}
\usetikzlibrary {arrows.meta}
\usetikzlibrary{calc,positioning}
\usepgfplotslibrary{fillbetween}
\usetikzlibrary{angles,quotes}
\usetikzlibrary{tikzmark}


\definecolor{PastelRed}{HTML}{FFC1CC}
\definecolor{PastelBlue}{HTML}{C2F0FF}
\definecolor{PastelBlue2}{HTML}{3182BD}

\definecolor{PastelPink}{HTML}{FFCBE1}
\definecolor{PastelGreen}{HTML}{D6E5BD}
\definecolor{PastelYellow}{HTML}{F9E1A8}
\definecolor{PastelBlue}{HTML}{BCD8EC}
\definecolor{PastelPurple}{HTML}{DCCCEC}
\definecolor{PastelOrange}{HTML}{FFDAB4}


\renewcommand{\footnotesize}{\tiny}

\newcommand{\mespor}{
  \ifcase\month
    \or janeiro
    \or fevereiro
    \or março
    \or abril
    \or maio
    \or junho
    \or julho
    \or agosto
    \or setembro
    \or outubro
    \or novembro
    \or dezembro
  \fi
}
% \usepackage{csquotes}
% \usepackage{minted}
% \usemintedstyle{xcode}
\definecolor{vscLightBackground}{HTML}{FFFFFF}
\definecolor{vscLightText}{HTML}{000000}
\definecolor{vscLightGreen}{HTML}{008000}        % For Comments AND Numbers
\definecolor{vscLightRed}{HTML}{A31515}          % For Strings
\definecolor{vscLightTeal}{HTML}{267F99}
\definecolor{vscBlue}{HTML}{0000FF}% For main keywords (if, for, def...)
\definecolor{vscLightPytorch}{HTML}{800080}      % For PyTorch keywords (dark magenta)

\lstdefinestyle{vscode-light}{
    language=Python,
    backgroundcolor=\color{vscLightBackground},
    basicstyle=\ttfamily\scriptsize\color{vscLightText},
    keywordstyle=\color{vscLightTeal},
    commentstyle=\color{vscLightGreen},
    stringstyle=\color{vscLightRed},
    % Style for PyTorch keywords (using keyword group 2)
    keywordstyle=[2]\color{vscLightPytorch}, 
    morekeywords={self, True, False, None},
    morekeywords=[2]{
        torch, Tensor, device, dtype, to, nn, optim, autograd, utils, data,
        Module, Linear, Conv2d, ReLU, CrossEntropyLoss, Adam, SGD,
        Dataset, DataLoader, randn, zeros, ones, arange, cat, stack,
        view, reshape, sum, mean, max, min, matmul, no_grad, backward
    },
    keywordstyle=[3]\color{vscBlue}, 
    morekeywords=[3]{
        model, loss, y, y_hat, x , optimizer
    },
    frame=single,
    tabsize=1,
    showstringspaces=false,
    breaklines=true,
    captionpos=b,
    extendedchars=true,
    numbers=left,
}

\usepackage[cache=true]{minted}
% minted v3+ (TeX Live 2024+) auto-detects the output directory.
% minted v2 (Ubuntu apt TeX Live 2023) needs it set explicitly to
% match latexmkrc's $out_dir = 'build'.
\makeatletter
\@ifpackagelater{minted}{2024/01/01}{}{%
  \def\minted@outputdir{build/}%
}
\makeatother
\setminted{
    fontsize=\scriptsize,
    style=vs,
    breaklines=true,
    numbers=left,
    numbersep=5pt,
    frame=single,
}
\providecommand{\citeauthoryear}[1]{(\citeauthor{#1},~\citeyear{#1})}

\providecommand{\thetav}{\bm{\theta}}

% vetor coluna 2D com colchetes
\providecommand{\cvec}[2]{\begin{bmatrix} #1 \\ #2 \end{bmatrix}}

\providecommand{\varvector}[1]{\mathbf{#1}}
% instance
\providecommand{\X}{\varvector{X}}
\providecommand{\mlinstance}{\varvector{X}}
\providecommand{\mlinstancei}{\varvector{X}^{(i)}}
\providecommand{\mlinstancex}[1]{\varvector{X}^{(#1)}}

% target
\providecommand{\y}{\varvector{y}}
\providecommand{\mltarget}{\varvector{y}}
\providecommand{\mltargeti}{\varvector{y}^{(i)}}
\providecommand{\mltargetx}[1]{\varvector{y}^{(#1)}}

%hats
\providecommand{\yhat}{\hat{\varvector{y}}}
\providecommand{\mltargethat}{\hat{\varvector{y}}}
\providecommand{\mltargethati}{\hat{\varvector{y}}^{(i)}}
\providecommand{\mltargethatx}[1]{\hat{\varvector{y}}^{(#1)}}

% linear reg
\providecommand{\mllinearregression}{\mltargethat=\mlinstance\bm{\theta}}

\providecommand{\tikzarrowcolor}[3]{%
  \begin{tikzpicture}[remember picture, overlay]
    % Arrow body
    \path[fill=#3] ([shift={(#1,#2)}]current page.south west) 
      rectangle ++(0.3, -0.2);
    % Arrow head
    \path[fill=#3] ([shift={(#1+0.3,#2+0.1)}]current page.south west) -- 
         ++(0, -0.4) -- 
         ++(0.2, 0.2) -- 
         cycle;
  \end{tikzpicture}%
}

\providecommand{\tikzarrow}[2]{%
    \tikzarrowcolor{#1}{#2}{Red}
}



\DeclareMathOperator*{\argmax}{argmax}
\DeclareMathOperator*{\argmin}{argmin}



% Bibliotecas TikZ adicionais usadas neste deck
\usetikzlibrary{matrix}

% Cores locais para consistência com o material original
\definecolor{qcolor}{RGB}{139,76,177}      % roxo (queries)
\definecolor{kcolor}{RGB}{218,165,32}      % dourado (keys)
\definecolor{vcolor}{RGB}{52,152,219}      % azul (values)
\definecolor{embcolor}{RGB}{120,180,90}    % verde (embeddings)
\definecolor{lightyellow}{RGB}{255,252,220}
\definecolor{lightblue}{RGB}{220,235,252}
\definecolor{lightgreen}{RGB}{220,245,220}
\definecolor{lightred}{RGB}{255,225,225}
\definecolor{lightpurple}{RGB}{230,220,255}
\definecolor{lightorange}{RGB}{255,237,210}
\definecolor{darkgreen}{RGB}{0,120,0}
\definecolor{encblock}{RGB}{215,238,210}
\definecolor{decblock}{RGB}{255,225,210}
\definecolor{mhablock}{RGB}{255,235,180}
\definecolor{ffblock}{RGB}{210,225,250}
\definecolor{normblock}{RGB}{250,235,200}

% Atalhos para destacar Q, K, V em fórmulas
\newcommand{\Qc}{\textcolor{qcolor}{Q}}
\newcommand{\Kc}{\textcolor{kcolor}{K}}
\newcommand{\Vc}{\textcolor{vcolor}{V}}

% Slide da arquitetura completa em duas colunas.
% #1: texto da coluna esquerda; #2: destaques desenhados sobre a figura,
% em coordenadas relativas ((0,0) = canto inferior esquerdo, (1,1) = superior direito).
\tikzset{foco/.style={draw=red, line width=1.6pt}}
\newcommand{\archslide}[2]{%
\begin{frame}{A Arquitetura Completa}
\begin{columns}[c]
\begin{column}{0.46\textwidth}
#1
\end{column}
\begin{column}{0.50\textwidth}
\centering
\begin{tikzpicture}
    \node[anchor=south west, inner sep=0] (img) at (0,0) {\includegraphics[height=0.80\textheight]{slides/img/04-transformers/vaswani2017-arquitetura.png}};
    \begin{scope}[x={(img.south east)}, y={(img.north west)}, overlay]
        #2
    \end{scope}
\end{tikzpicture}

\vspace{2pt}
{\footnotesize Figura de Vaswani et al., 2017. \textit{Attention is All You Need.}\par}
\end{column}
\end{columns}
\end{frame}}

% Título
\title{Deep Learning}
\subtitle{Transformers: Da Atenção em RNNs à Arquitetura Completa}
\author{Lucas Silveira Kupssinskü}
\institute{Machine Learning Theory and Applications Laboratory \\ PUCRS}
\date{\mespor de \the\year}

\begin{document}

% ============================================================
\begin{frame}
    \titlepage
\end{frame}

% ============================================================
\begin{frame}{Visão Geral}
    \tableofcontents
\end{frame}

% ============================================================
\section{Atenção em RNNs}
% ============================================================

\begin{frame}{O Problema do Gargalo (Encoder-Decoder)}

Em arquiteturas \textit{seq2seq} clássicas com RNNs, o \textbf{encoder} comprime toda a sequência de entrada em um único vetor de estado.

\vspace{6pt}

\begin{block}{O gargalo}
\begin{itemize}
    \item Frases longas precisam caber em um vetor de tamanho fixo
    \item O \textbf{decoder} acessa apenas o último estado oculto do encoder
    \item Informação inicial da sequência tende a se perder
\end{itemize}
\end{block}

\vspace{6pt}

\begin{exampleblock}{Solução: mecanismo de \textit{atenção}}
A \textbf{atenção} é um bloco da rede que conecta encoder e decoder e atribui \textit{scores} a cada etapa do encoder. Interpretamos esses scores como uma decisão de quais partes da sequência de entrada são mais relevantes para gerar o próximo token.
\end{exampleblock}

\vspace{2pt}
\centering
\footnotesize Bahdanau et al., 2014. \textit{Neural machine translation by jointly learning to align and translate.}
\end{frame}

% ============================================================

\begin{frame}{Encoder-Decoder com Atenção (Esquema)}

\begin{center}
\begin{tikzpicture}[
    >=Stealth, font=\small,
    enc/.style={draw, fill=encblock, minimum width=0.7cm, minimum height=1.0cm},
    dec/.style={draw, fill=decblock, minimum width=0.7cm, minimum height=1.0cm},
    tok/.style={font=\scriptsize, anchor=north},
]
    % Encoder cells
    \foreach \i/\w in {1/I, 2/saw, 3/cat, 4/on, 5/mat} {
        \node[enc] (e\i) at (\i*1.0, 0) {};
        \node[tok] at (e\i.south) {\w};
    }
    % EOS
    \node[enc] (e6) at (6.0, 0) {};
    \node[tok] at (e6.south) {\textless eos\textgreater};

    % Encoder arrows
    \foreach \i/\j in {1/2, 2/3, 3/4, 4/5, 5/6} {
        \draw[->, darkgreen!70] (e\i.east) -- (e\j.west);
    }

    % Encoder label
    \node[font=\small\bfseries, darkgreen] at (3.5, -1.4) {Encoder (RNN)};

    % Decoder cells
    \foreach \i/\w in {1/{\textless bos\textgreater}, 2/I, 3/saw, 4/a} {
        \node[dec] (d\i) at (8.0 + \i*1.0, 0) {};
        \node[tok] at (d\i.south) {\w};
    }

    % Decoder arrows
    \foreach \i/\j in {1/2, 2/3, 3/4} {
        \draw[->, red!70] (d\i.east) -- (d\j.west);
    }

    % Decoder label
    \node[font=\small\bfseries, red!70!black] at (10.5, -1.4) {Decoder};

    % Attention: from current decoder step to all encoder cells
    \foreach \i in {1,...,6} {
        \draw[->, gray, thin, opacity=0.7] (d3.north) to[bend left=10] (e\i.north);
    }

    % Context node
    \node[draw, fill=lightblue, rounded corners, minimum width=1.6cm, minimum height=0.5cm, font=\scriptsize] (ctx) at (3.5, 2.2) {contexto $c^{(t)}$};

    \draw[->, thick] (ctx.east) -- ++(2,0) -| (d3.north);

    \node[font=\scriptsize, anchor=west] at (ctx.east) {\hspace{2cm}};

    % Score arrows
    \foreach \i in {1,...,6} {
        \draw[->, gray, dotted, opacity=0.6] (e\i.north) -- (ctx.south);
    }

\end{tikzpicture}
\end{center}

\vspace{4pt}

A cada passo $t$ do decoder calculamos um \textbf{contexto} $c^{(t)}$, uma soma ponderada dos estados do encoder, com pesos dados pela atenção.

\end{frame}

% ============================================================

\begin{frame}[shrink=10]{Atenção: Scores, Pesos e Contexto}

A cada passo $t$ do decoder com estado $h_t$ e estados do encoder $s_1, \dots, s_m$:

\vspace{4pt}

\begin{block}{Pipeline}
\begin{enumerate}
    \item \textbf{Score} entre o estado atual e cada estado da fonte:
    \[
        \mathrm{score}(h_t, s_k), \quad k = 1, 2, \dots, m
    \]
    \item \textbf{Pesos da atenção} (softmax sobre os scores):
    \[
        a_k^{(t)} = \frac{\exp\bigl(\mathrm{score}(h_t, s_k)\bigr)}{\sum_{i=1}^{m} \exp\bigl(\mathrm{score}(h_t, s_i)\bigr)}
    \]
    \item \textbf{Contexto} (soma ponderada dos estados da fonte):
    \[
        c^{(t)} = \sum_{k=1}^{m} a_k^{(t)} \, s_k
    \]
\end{enumerate}
\end{block}

\vspace{4pt}

O contexto $c^{(t)}$ é então combinado ao estado do decoder para produzir a próxima saída.

\end{frame}

% ============================================================

\begin{frame}{Como Calcular os Scores}

Existem várias formas de definir $\mathrm{score}(h_t, s_k)$:

\vspace{6pt}

\begin{center}
\renewcommand{\arraystretch}{1.6}
\begin{tabular}{lcl}
\toprule
\textbf{Variante} & \textbf{Fórmula} & \textbf{Referência} \\
\midrule
Dot-product   & $\mathrm{score}(h_t, s_k) = h_t^{\top} s_k$ & --- \\
Bilinear      & $\mathrm{score}(h_t, s_k) = h_t^{\top} W \, s_k$ & Luong et al., 2015 \\
MLP (aditiva) & $\mathrm{score}(h_t, s_k) = w_2^{\top} \tanh\bigl(W_1 [h_t; s_k]\bigr)$ & Bahdanau et al., 2014 \\
\bottomrule
\end{tabular}
\end{center}

\vspace{8pt}

\begin{itemize}
    \item \textbf{Dot-product:} mais simples, sem parâmetros adicionais
    \item \textbf{Bilinear:} adiciona uma matriz de pesos $W$ aprendida
    \item \textbf{MLP:} mais expressiva, mas mais cara computacionalmente
\end{itemize}

\vspace{4pt}
\centering
A versão \textbf{dot-product} domina as arquiteturas modernas (transformers).

\end{frame}

% ============================================================
\section{Scaled Dot Product Attention}
% ============================================================

\begin{frame}{Por que ``scaled''?}

Treinar redes com \textit{dot product attention} pode causar \textbf{vanishing gradient} quando as dimensões dos vetores são grandes.

\vspace{6pt}

\begin{block}{Intuição}
Para vetores de dimensão $d$ grande, $h_t^{\top} s_k$ tende a ter magnitudes grandes. A softmax sobre esses scores satura, fazendo com que quase todos os pesos da atenção fiquem próximos de 0 ou de 1.
\end{block}

\vspace{6pt}

\textbf{Vamos quantificar:} suponha $s_i, h_i \sim \mathcal{N}(0, 1)$ i.i.d., componentes dos vetores $S, H \in \mathbb{R}^d$.

\vspace{4pt}
Calcule $\mathbb{E}[S \cdot H]$ e $\mathrm{var}[S \cdot H]$.

\end{frame}

% ============================================================

\begin{frame}{Esperança e Variância do Produto Escalar}

\begin{columns}[T]
\begin{column}{0.42\textwidth}
\textbf{Esperança:}
\begin{align*}
\mathbb{E}[S \cdot H]
&= \mathbb{E}\!\left[\sum_{i=1}^{d} s_i h_i\right] \\
&= \sum_{i=1}^{d} \mathbb{E}[s_i]\,\mathbb{E}[h_i] = 0
\end{align*}
\onslide<2->{%
\vspace{6pt}
{\footnotesize Para $s_i, h_i \sim \mathcal{N}(0, 1)$ i.i.d., $\mathbb{E}[s_i^2] = \mathrm{var}(s_i) = 1$ (pois $\mathbb{E}[s_i] = 0$), o que justifica os dois últimos passos.}

\vspace{6pt}
\textbf{Conclusão:} $\mathrm{var}[S \cdot H]$ cresce linearmente com $d$.
}
\end{column}
\begin{column}{0.58\textwidth}
\onslide<2->{%
\textbf{Variância:} por independência,
\begin{align*}
\mathrm{var}[S \cdot H]
&= \sum_{i=1}^{d} \mathrm{var}[s_i h_i] \\[2pt]
&= \sum_{i=1}^{d} \Bigl(\mathbb{E}[s_i^2 h_i^2] - \mathbb{E}[s_i h_i]^2\Bigr) \\[2pt]
&= \sum_{i=1}^{d} \mathbb{E}[s_i^2]\,\mathbb{E}[h_i^2] \\[2pt]
&= \sum_{i=1}^{d} \mathrm{var}(s_i)\,\mathrm{var}(h_i) = d
\end{align*}
}
\end{column}
\end{columns}

\end{frame}

% ============================================================

\begin{frame}{Scaled Dot Product Attention}

Para evitar \textit{scores} com magnitude crescente em $d$, normalizamos por $\sqrt{d}$:

\vspace{6pt}

\begin{exampleblock}{Scaled Dot Product Attention}
\[
\mathrm{score}(h_t, s_k) = \frac{h_t^{\top} s_k}{\sqrt{d}}, \quad k = 1, 2, \dots, m
\]
\[
a_k^{(t)} = \frac{\exp(\mathrm{score}(h_t, s_k))}{\sum_{i=1}^{m} \exp(\mathrm{score}(h_t, s_i))}, \qquad
c^{(t)} = \sum_{k=1}^{m} a_k^{(t)} s_k
\]
\end{exampleblock}

\vspace{6pt}

\begin{itemize}
    \item Escalar por $\sqrt{d}$ mantém a variância dos scores em $1$
    \item A softmax recebe valores em uma faixa estável: gradientes mais saudáveis
    \item Componente fundamental do transformer (Vaswani et al., 2017)
\end{itemize}

\end{frame}

% ============================================================
\section{Self-Attention}
% ============================================================

\begin{frame}{E se Aplicarmos Atenção da Sequência sobre Ela Mesma?}

Até aqui usamos atenção entre o estado do decoder e os estados do encoder (RNN).

\vspace{6pt}

\begin{block}{Ideia central de \textit{self-attention}}
E se utilizarmos \textit{Dot Product Attention} nos \textbf{tokens de entrada contra eles mesmos}?
\begin{itemize}
    \item Poderíamos codificar as relações entre tokens da entrada
    \item Não precisaríamos mais de RNNs como encoder
\end{itemize}
\end{block}

\vspace{8pt}

\begin{exampleblock}{Por exemplo, considere a frase}
\centering
\texttt{"The animal didn't cross the street because \textcolor{red}{it} was tired"}
\end{exampleblock}

\vspace{6pt}

A palavra \texttt{it} se refere a quais outras palavras da frase? Podemos usar \textit{attention} para descobrir.

\vspace{2pt}
\centering
\footnotesize Vaswani, A., Shazeer, N., Parmar, N., Uszkoreit, J., Jones, L., Gomez, A. N., Kaiser, {\L}., \& Polosukhin, I. (2017). Attention is All You Need. \textit{Advances in Neural Information Processing Systems}, v.\ 30.

\end{frame}

% ============================================================

\begin{frame}{Intuição: atenção como um dicionário ``suave''}

No tipo \textbf{dicionário}, cada \textbf{chave} (\textit{key}) tem um \textbf{valor} (\textit{value}), e uma \textbf{consulta} (\textit{query}) devolve o valor da chave correspondente.

\vspace{6pt}
\begin{columns}[T]
\begin{column}{0.5\textwidth}
\begin{block}{Dicionário comum}
\centering
\texttt{dic = \{1: 10, 2: 50, 3: 100\}}
\begin{itemize}
    \item \texttt{dic[2]} $\rightarrow$ \texttt{50} (chave exata)
    \item \texttt{dic[1.5]} $\rightarrow$ \textcolor{red}{\texttt{KeyError}}
\end{itemize}
\end{block}
\end{column}
\begin{column}{0.5\textwidth}
\onslide<2->{%
\begin{exampleblock}{Dicionário ``suave''}
Não dá erro: devolve uma \textbf{média dos valores ponderada pela proximidade} às chaves.
\begin{center}
\texttt{dic[1.5]} $\rightarrow$ \texttt{30} \quad {\scriptsize $\left(\tfrac{10+50}{2}\right)$}
\end{center}
\end{exampleblock}
}
\end{column}
\end{columns}

\vspace{6pt}
\onslide<2->{%
\begin{block}{É exatamente o \textit{self-attention}}
A \textcolor{qcolor}{\textbf{query}} é comparada a todas as \textcolor{kcolor}{\textbf{keys}}; quanto maior a similaridade (o \textit{score}), mais peso o \textcolor{vcolor}{\textbf{value}} correspondente recebe na média final.
\end{block}
}

\end{frame}

% ============================================================

\begin{frame}{Self-Attention: a Forma ``Crua'' Não Funciona}

Aplicar dot product attention diretamente sobre os tokens da entrada $x_1, \dots, x_m$:

\vspace{4pt}

\[
\mathrm{score}(x_t, x_k) = \frac{x_t^{\top} x_k}{\sqrt{d}}, \qquad
a_k^{(t)} = \mathrm{softmax}_k\bigl(\mathrm{score}(x_t, x_k)\bigr)
\]
\[
c^{(t)} = \sum_{k=1}^{m} a_k^{(t)} \, x_k
\]

\vspace{6pt}

\begin{alertblock}{Problema}
Não há \textbf{nada para aprender} aqui: o cálculo é determinístico nas entradas. Sem parâmetros treináveis, \textit{backpropagation} não tem o que ajustar.
\end{alertblock}

\vspace{6pt}

\textbf{Solução:} introduzir três projeções lineares aprendidas para cada token, gerando \textbf{queries}, \textbf{keys} e \textbf{values}.

\end{frame}

% ============================================================

\begin{frame}[shrink=10]{Queries, Keys e Values}

Para cada token $x_t \in \mathbb{R}^{d_{\text{emb}}}$, projetamos em três vetores via matrizes aprendidas:

\vspace{6pt}

\begin{columns}[T]
\begin{column}{0.52\textwidth}
\begin{flalign*}
\textcolor{qcolor}{q_t} &= x_t \, W_q + b_q, & W_q &\in \mathbb{R}^{d_{\text{emb}} \times d}, & q_t &\in \mathbb{R}^{d} && \\
\textcolor{kcolor}{k_t} &= x_t \, W_k + b_k, & W_k &\in \mathbb{R}^{d_{\text{emb}} \times d}, & k_t &\in \mathbb{R}^{d} && \\
\textcolor{vcolor}{v_t} &= x_t \, W_v + b_v, & W_v &\in \mathbb{R}^{d_{\text{emb}} \times d_v}, & v_t &\in \mathbb{R}^{d_v} &&
\end{flalign*}
\end{column}
\begin{column}{0.48\textwidth}
\textbf{Analogia (recuperação de informação):}
\begin{itemize}
    \item \textcolor{qcolor}{\textbf{Query}}: vetor que codifica uma consulta
    \item \textcolor{kcolor}{\textbf{Key}}: chaves de busca, comparadas às queries
    \item \textcolor{vcolor}{\textbf{Value}}: valor retornado quando query e key são similares
\end{itemize}
\end{column}
\end{columns}

\vspace{6pt}

A atenção passa a ser:
\[
\mathrm{score}(q_t, k_i) = \frac{q_t^{\top} k_i}{\sqrt{d}}, \qquad
a_i^{(t)} = \mathrm{softmax}_i(\mathrm{score}(q_t, k_i))
\]
\[
c^{(t)} = \sum_{i=1}^{m} a_i^{(t)} \, v_i
\]

\textbf{Agora há parâmetros aprendidos:} $W_q, W_k, W_v$.

\end{frame}

% ============================================================

\begin{frame}[shrink=10]{Self-Attention: Exemplo Numérico (Setup)}

\textbf{Embeddings de entrada} (dois tokens, $d_{\text{emb}}=3$):
\[
x_1 = [1,\ -3,\ 2], \qquad x_2 = [2,\ -1,\ 1]
\]

\textbf{Matrizes de pesos} (todas com bias $0$, projeção para $d=2$):
\[
W_q = \begin{bmatrix} 1 & 0 \\ 1 & -1 \\ 0 & -2 \end{bmatrix}, \quad
W_k = \begin{bmatrix} 1 & 1 \\ 0 & -1 \\ 0 & 0 \end{bmatrix}, \quad
W_v = \begin{bmatrix} -2 & 1 \\ 0 & 1 \\ 0 & 2 \end{bmatrix}
\]

\vspace{4pt}
\centering
Projetamos cada token aplicando $W_q, W_k, W_v$; os cálculos completos vêm a seguir.

\end{frame}

% ============================================================

\begin{frame}[shrink=10]{Self-Attention: Cálculo das Projeções}

Cada token ($x_1 = [1,-3,2]$, $x_2 = [2,-1,1]$) é projetado por $W_q, W_k, W_v$:

\vspace{2pt}
\textcolor{qcolor}{\textbf{Queries}} ($q_i = x_i W_q$):
\begin{align*}
q_1 &= [\,1(1){+}(-3)(1){+}2(0),\ \ 1(0){+}(-3)(-1){+}2(-2)\,] = [-2,\ -1] \\
q_2 &= [\,2(1){+}(-1)(1){+}1(0),\ \ 2(0){+}(-1)(-1){+}1(-2)\,] = [1,\ -1]
\end{align*}

\textcolor{kcolor}{\textbf{Keys}} ($k_i = x_i W_k$):
\begin{align*}
k_1 &= [\,1(1){+}(-3)(0){+}2(0),\ \ 1(1){+}(-3)(-1){+}2(0)\,] = [1,\ 4] \\
k_2 &= [\,2(1){+}(-1)(0){+}1(0),\ \ 2(1){+}(-1)(-1){+}1(0)\,] = [2,\ 3]
\end{align*}

\textcolor{vcolor}{\textbf{Values}} ($v_i = x_i W_v$):
\begin{align*}
v_1 &= [\,1(-2){+}(-3)(0){+}2(0),\ \ 1(1){+}(-3)(1){+}2(2)\,] = [-2,\ 2] \\
v_2 &= [\,2(-2){+}(-1)(0){+}1(0),\ \ 2(1){+}(-1)(1){+}1(2)\,] = [-4,\ 3]
\end{align*}

\end{frame}

% ============================================================

\begin{frame}[shrink=15]{Self-Attention: Exemplo Numérico (Cálculo)}

\textbf{Foco no token 1:} usamos $\textcolor{qcolor}{q_1 = [-2, -1]}$ contra todas as keys.

\vspace{6pt}

\textbf{Scores não normalizados:}
\[
q_1 \cdot k_1 = (-2)(1) + (-1)(4) = -6, \qquad
q_1 \cdot k_2 = (-2)(2) + (-1)(3) = -7
\]

\vspace{4pt}

\textbf{Normalização por $\sqrt{d}$ (com $d=2$):}
\[
-6 / \sqrt{2} \approx -4.24, \qquad -7 / \sqrt{2} \approx -4.95
\]

\vspace{4pt}

\textbf{Softmax sobre os scores normalizados:}
\[
a_1^{(1)} \approx 0.67, \qquad a_2^{(1)} \approx 0.33
\]

\vspace{4pt}

\textbf{Soma ponderada dos values:}
\[
c^{(1)} = 0.67 \cdot [-2, 2] + 0.33 \cdot [-4, 3] \approx [-2.66,\ 2.33]
\]

\vspace{4pt}
\centering
O token 1 atende \textbf{principalmente} a si mesmo (0.67), mas incorpora uma fração do token 2 (0.33): a saída é uma média ponderada.

\end{frame}

% ============================================================

\begin{frame}{Forma Matricial: Processar Toda a Sequência de uma Vez}

A implementação anterior processa cada query \textbf{de forma independente}. Como o resultado de cada query não depende das demais (no mesmo passo), podemos vetorizar tudo:

\vspace{4pt}

\begin{exampleblock}{Scaled Dot Product Attention (forma matricial)}
\[
\boxed{\;
\mathrm{attention}(\Qc, \Kc, \Vc) = \mathrm{softmax}\!\left(\frac{\Qc \, \Kc^{\top}}{\sqrt{d}}\right) \Vc
\;}
\]
\end{exampleblock}

\vspace{6pt}

\textbf{Com} ($X$: \texttt{seq\_size $\times$ emb\_size}):
\begin{itemize}
    \item $\Qc = X W_q$, $\quad$ $W_q$: \texttt{emb\_size $\times$ d}, $\quad$ $\Qc$: \texttt{seq\_size $\times$ d}
    \item $\Kc = X W_k$, $\quad$ $W_k$: \texttt{emb\_size $\times$ d}, $\quad$ $\Kc$: \texttt{seq\_size $\times$ d}
    \item $\Vc = X W_v$, $\quad$ $W_v$: \texttt{emb\_size $\times$ value\_size}, $\quad$ $\Vc$: \texttt{seq\_size $\times$ value\_size}
\end{itemize}

\vspace{4pt}

A matriz $\mathrm{softmax}(\Qc \Kc^{\top} / \sqrt{d})$ tem shape \texttt{seq\_size $\times$ seq\_size}: contém os \textbf{pesos da atenção} de cada token sobre cada token.

\end{frame}

% ============================================================

\begin{frame}[shrink=10]{Forma Matricial: Projeções $\Qc, \Kc, \Vc$}

Empilhando os dois embeddings nas linhas de $X$ ($2\times3$) e projetando com $W_q, W_k, W_v$ ($3\times2$), cada projeção fica $2\times2$:

\vspace{2pt}
\[
\Qc = X W_q =
\begin{bmatrix} 1 & -3 & 2 \\ 2 & -1 & 1 \end{bmatrix}
\begin{bmatrix} 1 & 0 \\ 1 & -1 \\ 0 & -2 \end{bmatrix}
= \begin{bmatrix} -2 & -1 \\ 1 & -1 \end{bmatrix}
\]
\[
\Kc = X W_k =
\begin{bmatrix} 1 & -3 & 2 \\ 2 & -1 & 1 \end{bmatrix}
\begin{bmatrix} 1 & 1 \\ 0 & -1 \\ 0 & 0 \end{bmatrix}
= \begin{bmatrix} 1 & 4 \\ 2 & 3 \end{bmatrix}
\]
\[
\Vc = X W_v =
\begin{bmatrix} 1 & -3 & 2 \\ 2 & -1 & 1 \end{bmatrix}
\begin{bmatrix} -2 & 1 \\ 0 & 1 \\ 0 & 2 \end{bmatrix}
= \begin{bmatrix} -2 & 2 \\ -4 & 3 \end{bmatrix}
\]

\vspace{2pt}
\centering
Cada \textbf{linha} de $\Qc, \Kc, \Vc$ é o vetor projetado de um token.

\end{frame}

% ============================================================

\begin{frame}[shrink=10]{Forma Matricial: Pipeline Completo}

Com as projeções do slide anterior:
\[
\Qc = \begin{bmatrix} -2 & -1 \\ 1 & -1 \end{bmatrix}, \quad
\Kc = \begin{bmatrix} 1 & 4 \\ 2 & 3 \end{bmatrix}, \quad
\Vc = \begin{bmatrix} -2 & 2 \\ -4 & 3 \end{bmatrix}
\]

\vspace{6pt}

\textbf{Scores e normalização:}
\[
\Qc \, \Kc^{\top} = \begin{bmatrix} -6 & -7 \\ -3 & -1 \end{bmatrix}, \qquad
\frac{\Qc \, \Kc^{\top}}{\sqrt{2}} \approx \begin{bmatrix} -4.24 & -4.95 \\ -2.12 & -0.71 \end{bmatrix}
\]

\vspace{4pt}

\textbf{Softmax (por linha):}
\[
\mathrm{softmax} \approx \begin{bmatrix} 0.67 & 0.33 \\ 0.20 & 0.80 \end{bmatrix}
\]

\vspace{4pt}

\textbf{Saída final:}
\[
\mathrm{attention}(\Qc, \Kc, \Vc) \approx \begin{bmatrix} -2.66 & 2.33 \\ -3.61 & 2.80 \end{bmatrix}
\]

\end{frame}

% ============================================================

\begin{frame}{Refletindo sobre Self-Attention}

\begin{block}{É não linear?}
\textbf{Sim:} $Q K^{\top} V$ é polinomial de grau 3 nos parâmetros; a softmax já é não-linear.
\end{block}

\begin{columns}[T]
\column{0.48\textwidth}
\begin{block}{Vantagem em relação a um MLP}
O número de parâmetros \textbf{não escala} com o tamanho da sequência. Os mesmos $W_q, W_k, W_v$ atendem qualquer comprimento.
\end{block}

\column{0.48\textwidth}
\begin{block}{Vantagem em relação a RNNs}
\begin{itemize}
    \item Muito mais fácil de \textbf{paralelizar}: tudo é matricial
    \item Lida melhor com \textbf{sequências longas}: dependências de longo alcance em um único passo
\end{itemize}
\end{block}
\end{columns}

\vspace{3pt}

\begin{exampleblock}{Bônus: explicabilidade}
Os pesos da atenção formam uma matriz interpretável: dá para plotar e ver quais tokens a rede ``olha''.
\end{exampleblock}

\end{frame}

% ============================================================

\begin{frame}[shrink=10]{O Calcanhar de Aquiles: Falta Informação Posicional}

\begin{alertblock}{Problema}
Da forma como definimos, \textit{self-attention} \textbf{não tem informação sobre a posição} dos tokens na sequência.
\end{alertblock}

\vspace{6pt}

Em RNNs e CNNs, a informação posicional está integrada ao \textit{viés indutivo} dos modelos, então não demandavam atenção especial. Em self-attention, não.

\vspace{8pt}

\begin{exampleblock}{Consequência absurda}
\centering
\texttt{"lucas comeu pastel no jantar"} $=$ \texttt{"jantar comeu pastel no lucas"}

\vspace{4pt}

Para o modelo, ambas as frases produzem o mesmo conjunto de pesos da atenção.
\end{exampleblock}

\vspace{6pt}

\centering
Precisamos adaptar os embeddings para que carreguem informação sobre a \textbf{posição do token} na sequência.

\end{frame}

% ============================================================
\section{Positional Encoding}
% ============================================================

\begin{frame}{Ideia Geral do Positional Encoding}

\textbf{Ideia:} criar uma codificação para a posição e \textbf{somá-la} ao embedding do token.

\vspace{8pt}

\begin{center}
\begin{tikzpicture}[
    >=Stealth, font=\small,
    box/.style={draw, minimum width=1.6cm, minimum height=1.0cm},
]
    \node[box, fill=lightpurple] (we) at (0, 0) {word emb.};
    \node[font=\Large] at (2.0, 0) {$+$};
    \node[box, fill=lightyellow] (pe) at (4.0, 0) {position emb.};
    \node[font=\Large] at (6.0, 0) {$=$};
    \node[box, fill=lightred] (out) at (8.4, 0) {input do transformer};

    \node[font=\scriptsize, anchor=north] at (we.south) {tamanho $d$};
    \node[font=\scriptsize, anchor=north] at (pe.south) {tamanho $d$};
    \node[font=\scriptsize, anchor=north] at (out.south) {tamanho $d$};
\end{tikzpicture}
\end{center}

\vspace{6pt}

\begin{itemize}
    \item Existem várias formas de definir/aprender embeddings posicionais
    \item No paper original, Vaswani et al.\ (2017) observaram que \textbf{embeddings posicionais aprendidos não melhoraram} resultados em relação a embeddings posicionais \textbf{fixos} (sinusoidal)
\end{itemize}

\end{frame}

% ============================================================

\begin{frame}{Propriedades Desejadas}

Precisamos de uma matriz de posições que codifique adequadamente posições \textbf{absolutas} e \textbf{relativas} dos tokens em uma sequência.

\vspace{8pt}

\begin{block}{Requisitos}
\begin{itemize}
    \item \textbf{Único por posição:} cada posição da sequência tem uma codificação distinta
    \item \textbf{Distâncias relativas consistentes:} o deslocamento entre duas posições é codificado de forma consistente, independentemente do tamanho da sequência
    \item \textbf{Generaliza para qualquer tamanho:} a codificação tem limite superior e inferior bem definidos
    \item \textbf{Determinística:} sem aleatoriedade na inferência
\end{itemize}
\end{block}

\end{frame}

% ============================================================

\begin{frame}{Ideia 1: Usar o Índice da Posição}

\textbf{Proposta:} codificar a posição $i$ pelo próprio número $i$ (ou $i / N$).

\vspace{6pt}

\begin{alertblock}{Problemas}
\begin{itemize}
    \item Sem limite superior bem definido (para sequências arbitrárias)
    \item Os valores ficam grandes para sequências longas, dominando o embedding do token
    \item Normalizar por $N$ resolve o limite, mas a mesma posição relativa passa a depender do tamanho da sequência
\end{itemize}
\end{alertblock}

\vspace{6pt}

\centering
Não funciona.

\end{frame}

% ============================================================

\begin{frame}{Ideia 2: Usar Números Binários}

\textbf{Proposta:} usar a representação binária da posição em $\lceil \log_2 N \rceil$ bits.

\vspace{4pt}

\begin{columns}[T]
\column{0.55\textwidth}
\begin{itemize}
    \item Existe um \textbf{comportamento periódico} natural nos bits binários
    \item Bits menos significativos alternam rápido (período 2), mais significativos alternam devagar
    \item Isso é desejável para codificar posições
\end{itemize}

\vspace{6pt}

\begin{alertblock}{Problema}
Redes neurais usam matemática \textbf{contínua}. Interpolar entre representações binárias não preserva a noção de distância. Também não funciona.
\end{alertblock}

\column{0.40\textwidth}
\centering
\footnotesize
\renewcommand{\arraystretch}{1.0}
\begin{tabular}{r|cccc}
$0$ & 0 & 0 & 0 & 0 \\
$1$ & 0 & 0 & 0 & 1 \\
$2$ & 0 & 0 & 1 & 0 \\
$3$ & 0 & 0 & 1 & 1 \\
$4$ & 0 & 1 & 0 & 0 \\
$5$ & 0 & 1 & 0 & 1 \\
$6$ & 0 & 1 & 1 & 0 \\
$7$ & 0 & 1 & 1 & 1 \\
$8$ & 1 & 0 & 0 & 0 \\
$\vdots$ &  &  &  &  \\
$15$ & 1 & 1 & 1 & 1 \\
\end{tabular}
\end{columns}

\end{frame}

% ============================================================

\begin{frame}{Ideia 3: Senos e Cossenos de Frequências Variadas}

\textbf{Proposta:} usar uma combinação de senos e cossenos de diferentes frequências para indicar a posição.

\vspace{2pt}

\begin{block}{Por que funciona}
\begin{itemize}
    \item \textbf{Limite máximo e mínimo bem definidos:} $\sin, \cos \in [-1, 1]$
    \item \textbf{Periodicidade multiescala:} cada dimensão tem sua própria frequência, gerando códigos únicos
    \item \textbf{Linearidade da rotação:} a posição (ângulo) pode ser modificada por uma transformação linear
\end{itemize}
\end{block}

\vspace{2pt}

\[
\begin{pmatrix} \cos(\theta + \phi) \\ \sin(\theta + \phi) \end{pmatrix}
=
\begin{pmatrix} \cos\phi & -\sin\phi \\ \sin\phi & \cos\phi \end{pmatrix}
\begin{pmatrix} \cos\theta \\ \sin\theta \end{pmatrix}
\]

\vspace{2pt}

Isso significa que o modelo consegue ``rotacionar'' as posições por uma matriz de pesos aprendida, capturando deslocamentos relativos.

\end{frame}

% ============================================================

\begin{frame}{Positional Encoding Sinusoidal: Fórmula}

\begin{exampleblock}{Definição (Vaswani et al., 2017)}
\[
P_{i,j} =
\begin{cases}
\sin\!\left(\dfrac{i}{n^{(2k/d)}}\right), & \text{se } j = 2k \\[12pt]
\cos\!\left(\dfrac{i}{n^{(2k/d)}}\right), & \text{se } j = 2k + 1
\end{cases}
\]
\end{exampleblock}

\vspace{4pt}

\textbf{Onde:}
\begin{itemize}
    \item $i$: índice da posição na sequência
    \item $j$: índice da dimensão do embedding ($j = 0, \dots, d-1$)
    \item $d$: dimensão do embedding
    \item $n$: hiperparâmetro de escala (no paper, $n = 10000$)
\end{itemize}

\vspace{4pt}

Cada par de dimensões $(2k, 2k+1)$ usa uma frequência diferente, gerando uma assinatura única para cada posição.

\end{frame}

% ============================================================

\begin{frame}[shrink=10]{Positional Encoding: Comportamento do Produto Escalar}

Uma propriedade interessante: o produto escalar entre dois positional encodings tem um pico claro na posição central e decai suavemente.

\vspace{4pt}

\begin{center}
\begin{tikzpicture}
\begin{axis}[
    width=10cm, height=5cm,
    xlabel={Posição},
    ylabel={\scriptsize $\frac{\text{dot product}}{\max(\text{dot product})}$},
    domain=0:256, samples=400,
    ymin=0.3, ymax=1.05,
    xmin=0, xmax=256,
    grid=both, grid style={gray!20},
    tick label style={font=\scriptsize},
    label style={font=\scriptsize},
]
% Curva real: produto escalar de PEs sinusoidais (d=128, n=10000) entre a
% posição 128 e todas as demais, normalizado pelo valor em 128.
\addplot[thick, blue] coordinates {
(0,0.333) (4,0.415) (8,0.422) (12,0.460) (16,0.363) (20,0.426) (24,0.423) (28,0.477) (32,0.384) (36,0.430) (40,0.501) (44,0.402) (48,0.452) (52,0.501) (56,0.417) (60,0.500) (64,0.477) (68,0.469) (72,0.509) (76,0.498) (80,0.495) (84,0.542) (88,0.509) (92,0.548) (96,0.557) (100,0.559) (104,0.589) (108,0.608) (112,0.620) (116,0.662) (120,0.714) (124,0.759) (128,1.000) (132,0.759) (136,0.714) (140,0.662) (144,0.620) (148,0.608) (152,0.589) (156,0.559) (160,0.557) (164,0.548) (168,0.509) (172,0.542) (176,0.495) (180,0.498) (184,0.509) (188,0.469) (192,0.477) (196,0.500) (200,0.417) (204,0.501) (208,0.452) (212,0.402) (216,0.501) (220,0.430) (224,0.384) (228,0.477) (232,0.423) (236,0.426) (240,0.363) (244,0.460) (248,0.422) (252,0.415) (256,0.333)
};
\end{axis}
\end{tikzpicture}
\end{center}

\vspace{4pt}

\begin{itemize}
    \item \textbf{Pico em 1} quando comparamos uma posição com ela mesma
    \item \textbf{Decaimento suave} para posições próximas
    \item \textbf{Padrão oscilatório} característico das senoides combinadas
\end{itemize}

Isso permite ao modelo perceber proximidade relativa entre tokens.

\end{frame}

% ============================================================
\section{Arquitetura Transformer}
% ============================================================

\archslide{
Finalmente estamos prontos para abrir o transformer. Na prática, já vimos quase todos os componentes.

\vspace{8pt}

\begin{itemize}
    \item À esquerda, o \textbf{encoder}: processa a sequência de entrada
    \item À direita, o \textbf{decoder}: gera a saída token a token
\end{itemize}
}{}

% ============================================================

\archslide{
\textbf{Em foco: Input Embedding.}

\vspace{6pt}

A porta de entrada do transformer: transforma índices de tokens em vetores.
}{
\draw[foco] (0.321,0.1675) ellipse (0.155 and 0.05);
}

% ============================================================

\begin{frame}{Input Embedding: uma Matriz Aprendida}

Após a tokenização, cada token da sequência é um \textbf{índice inteiro} do vocabulário $V$.

\vspace{6pt}

\begin{block}{A camada de embedding}
É uma \textbf{matriz de parâmetros aprendidos}
\[
E \in \mathbb{R}^{|V| \times d},
\]
com uma linha por token do vocabulário: a linha $t$ é o embedding do token de índice $t$. A matriz é treinada por backpropagation como qualquer outro peso da rede.
\end{block}

\vspace{6pt}

\begin{exampleblock}{Ordem de grandeza}
Com $|V| = 50\,000$ e $d = 512$, a camada de embedding tem $|V| \times d = 25{,}6$ milhões de parâmetros.
\end{exampleblock}

\end{frame}

% ============================================================

\begin{frame}{Selecionando o Embedding de um Token}

Como o modelo seleciona o embedding do token de índice $t$? O índice vira um vetor \textbf{one-hot} $u \in \{0,1\}^{|V|}$, com $1$ apenas na posição $t$, e o produto $u^{\top} E$ devolve exatamente a \textbf{linha $t$} de $E$.

\vspace{4pt}

\begin{exampleblock}{Exemplo com $|V| = 4$ e $d = 3$, token de índice $3$}
\[
u^{\top} E =
\begin{bmatrix} 0 & 0 & 1 & 0 \end{bmatrix}
\begin{bmatrix}
0.2 & -1.1 & 0.7 \\
1.5 & 0.3 & -0.4 \\
-0.8 & 0.9 & 2.1 \\
0.6 & -0.5 & 1.3
\end{bmatrix}
=
\begin{bmatrix} -0.8 & 0.9 & 2.1 \end{bmatrix}
= E_{3,:}
\]
\end{exampleblock}

\vspace{6pt}

Na prática, os frameworks implementam a operação como \textit{lookup}: indexam a linha $t$ diretamente, sem multiplicação de matrizes.

\end{frame}

% ============================================================

\archslide{
\textbf{Em foco: Multi-Head Attention.}

\vspace{6pt}

O bloco de atenção que já conhecemos, agora dividido em várias ``cabeças'' paralelas.
}{
\draw[foco] (0.322,0.3725) ellipse (0.16 and 0.055);
}

% ============================================================

\begin{frame}[shrink=10]{Multi-Head Attention}

\begin{columns}[T]
\column{0.58\textwidth}
\textbf{Ideia:} separar a computação do self-attention em $H$ ``cabeças'' paralelas.

\vspace{6pt}

\begin{block}{Por exemplo, com $D = 512$ e $H = 8$}
\begin{itemize}
    \item Cada cabeça trabalha com dimensão $D/H = 64$
    \item As $H$ cabeças computam atenção em paralelo
    \item As saídas são \textbf{concatenadas}
    \item Aplicamos uma projeção linear final
\end{itemize}
\end{block}

\vspace{4pt}

\begin{exampleblock}{Fórmula}
\small
\[
\mathrm{MultiHead}(\Qc, \Kc, \Vc) = \mathrm{Concat}(\mathrm{head}_1, \dots, \mathrm{head}_H) W_O
\]
\[
\mathrm{head}_i = \mathrm{attention}(\Qc W_Q^i, \Kc W_K^i, \Vc W_V^i)
\]
\end{exampleblock}

\column{0.40\textwidth}
\begin{center}
\begin{tikzpicture}[
    >=Stealth, font=\scriptsize,
    box/.style={draw, rounded corners, minimum width=1.1cm, minimum height=0.5cm, align=center},
]
    \node[box, fill=lightblue] (lq) at (-1.6, 0) {Linear};
    \node[box, fill=lightblue] (lk) at (0.0,  0) {Linear};
    \node[box, fill=lightblue] (lv) at (1.6,  0) {Linear};

    \node[font=\scriptsize, anchor=north] at (lq.south) {Q};
    \node[font=\scriptsize, anchor=north] at (lk.south) {K};
    \node[font=\scriptsize, anchor=north] at (lv.south) {V};

    \node[box, fill=lightpurple, minimum width=4.6cm] (sdpa) at (0, 1.1) {Scaled Dot Product Attn (h vezes)};
    \node[box, fill=lightyellow, minimum width=4.6cm] (cc) at (0, 2.0) {Concat};
    \node[box, fill=lightblue,   minimum width=4.6cm] (lo) at (0, 2.9) {Linear};

    \draw[->] (lq) -- (sdpa.south -| lq);
    \draw[->] (lk) -- (sdpa.south -| lk);
    \draw[->] (lv) -- (sdpa.south -| lv);
    \draw[->] (sdpa) -- (cc);
    \draw[->] (cc) -- (lo);
\end{tikzpicture}
\end{center}
\end{columns}

\end{frame}

% ============================================================

\archslide{
\textbf{Em foco: Add \& Norm.}

\vspace{6pt}

Skip connections e normalização em torno de cada sub-bloco do encoder e do decoder.
}{
\draw[foco] (0.322,0.5925) ellipse (0.15 and 0.035);
\draw[foco] (0.322,0.4275) ellipse (0.15 and 0.035);
}

% ============================================================

\begin{frame}{Skip Connections (Residuais)}

Note no diagrama do transformer que existe uma seta passando \textbf{ao redor} do bloco \textit{Multi-Head Attention} (e do bloco Feed Forward). Isso é uma \textbf{skip connection}.

\vspace{6pt}

\begin{exampleblock}{Add \& Norm}
\centering
\[
\text{output} = \mathrm{LayerNorm}\bigl(x + \mathrm{Subblock}(x)\bigr)
\]
\end{exampleblock}

\vspace{6pt}

\begin{itemize}
    \item Mantém o gradiente fluindo (mesma motivação da ResNet)
    \item Permite empilhar muitos blocos sem que o sinal se degrade
    \item ``\textit{Add}'' é a soma residual; ``\textit{Norm}'' é a normalização (LayerNorm)
\end{itemize}

\vspace{4pt}
\centering
Sem skip connections, o transformer não treinaria em profundidade.

\end{frame}

% ============================================================

\begin{frame}{Layer Normalization}

\begin{itemize}
    \item Similar ao \textbf{batch norm} visto anteriormente (ResNet), mas as estatísticas são computadas \textbf{por token} (não por mini-batch)
    \item Introduz dois parâmetros aprendidos por dimensão (escala e deslocamento)
\end{itemize}

\vspace{8pt}
\begin{exampleblock}{Forma fechada}
Para um \textit{token} $x_i \in \mathbb{R}^{d}$,
\[
    \mathrm{LN}(x_i) = \gamma \odot \frac{x_i - \mu_i}{\sigma_i} + \beta, \qquad
    \mu_i = \frac{1}{d}\sum_{k=1}^{d} x_{i,k}, \quad
    \sigma_i^2 = \frac{1}{d}\sum_{k=1}^{d}(x_{i,k} - \mu_i)^2.
\]
\end{exampleblock}

\vspace{8pt}
\centering
\footnotesize Layer Norm é mais adequada para sequências de comprimento variável.

\end{frame}

% ============================================================

\begin{frame}{Layer Norm vs. Batch Norm}

\begin{center}
\begin{tikzpicture}[font=\scriptsize, >=Stealth]
    % Batch Norm panel
    \node[font=\small\bfseries] at (0, 2.4) {Batch Norm};
    \draw[red!70, dashed, thick, rounded corners] (-1.6, 1.0) rectangle (-0.7, 2.0);
    \draw[red!70, dashed, thick, rounded corners] (-1.6, -0.4) rectangle (-0.7, 0.6);

    \node at (-1.15, 1.5) {col};
    \node at (-1.15, 0.1) {col};
    \node[anchor=west, font=\scriptsize, align=left] at (-0.5, 0.8) {normaliza\\por coluna\\(feature)};

    % Layer Norm panel
    \node[font=\small\bfseries] at (5.0, 2.4) {Layer Norm};
    \draw[darkgreen, dashed, thick, rounded corners] (3.4, 1.5) rectangle (5.6, 2.0);
    \draw[darkgreen, dashed, thick, rounded corners] (3.4, 0.5) rectangle (5.6, 1.0);
    \draw[darkgreen, dashed, thick, rounded corners] (3.4, -0.5) rectangle (5.6, 0.0);

    \node at (4.5, 1.75) {linha (token)};
    \node at (4.5, 0.75) {linha (token)};
    \node at (4.5, -0.25) {linha (token)};
    \node[anchor=west, font=\scriptsize, align=left] at (5.8, 0.75) {normaliza\\por linha\\(por token)};
\end{tikzpicture}
\end{center}

\vspace{6pt}
\begin{itemize}
    \item \textbf{Batch Norm}: normaliza cada \textbf{coluna} (feature) com estatísticas do mini-batch.
    \item \textbf{Layer Norm}: normaliza cada \textbf{linha} (token) com as estatísticas das próprias features, independente do batch.
\end{itemize}

\end{frame}

% ============================================================

\archslide{
\textbf{Em foco: Feed Forward.}

\vspace{6pt}

Um pequeno MLP aplicado a cada posição da sequência, de forma independente.
}{
\draw[foco] (0.322,0.5375) ellipse (0.15 and 0.05);
}

% ============================================================

\begin{frame}{Feed-Forward Network}

Após o bloco \textit{Multi-Head Attention} e o \textit{Add \& Norm}, o transformer aplica um pequeno MLP a cada posição.

\vspace{6pt}

\begin{exampleblock}{Estrutura}
\[
\mathrm{FFN}(x) = \mathrm{ReLU}(x \, W_1 + b_1) \, W_2 + b_2
\]
\end{exampleblock}

\vspace{6pt}

\begin{itemize}
    \item No paper original, são duas camadas com ReLU no meio
    \item A entrada do MLP \textbf{depende do tamanho da sequência}? Não: aplicado posição a posição (\textit{pointwise})
    \item Após interagir com os outros tokens via atenção, cada posição ``pensa'' por si mesma
    \item Geralmente expande a dimensão (ex.: $d \to 4d \to d$)
\end{itemize}

\end{frame}

% ============================================================

\archslide{
\textbf{Em foco: o bloco do encoder.}

\vspace{6pt}

Os componentes que acabamos de ver formam um bloco, replicado $N$ vezes.
}{
\draw[foco] (0.27,0.457) ellipse (0.245 and 0.21);
}

% ============================================================

\begin{frame}{O Bloco do Encoder e a Pilha $N\times$}

Cada \textbf{bloco do encoder} é a sequência:

\vspace{4pt}
\begin{center}
Multi-Head Attention $\to$ Add \& Norm $\to$ Feed Forward $\to$ Add \& Norm
\end{center}

\vspace{6pt}

\begin{block}{Empilhamento}
Replicamos esse bloco $N$ vezes (no paper original, $N = 6$). Cada bloco refina a representação dos tokens, integrando informação dos demais via self-attention.
\end{block}

\vspace{6pt}

\begin{itemize}
    \item A entrada do encoder é \texttt{Input Embedding $+$ Positional Encoding}
    \item A saída do encoder é uma sequência de vetores que, conjuntamente, codificam o significado da entrada \textbf{contextualizada}
    \item Essa saída será consumida pelo decoder via \textbf{cross-attention}
\end{itemize}

\end{frame}

% ============================================================

\archslide{
\textbf{Em foco: Masked Multi-Head Attention.}

\vspace{6pt}

O primeiro bloco de atenção do decoder só pode olhar para tokens já gerados.
}{
\draw[foco] (0.676,0.3825) ellipse (0.16 and 0.068);
}

% ============================================================

\begin{frame}{O Decoder: Masked Self-Attention}

O primeiro bloco de atenção do decoder é praticamente igual ao do encoder, com uma diferença crítica: aplica uma \textbf{máscara} para olhar apenas tokens anteriores.

\vspace{6pt}

\begin{block}{Por que mascarar?}
Durante o treinamento (tarefa de tradução, por exemplo), o decoder vê toda a sequência de saída. Sem máscara, ele \textit{trapacearia}, atendendo a tokens futuros que ainda não deveriam ter sido gerados. A máscara garante que cada posição $t$ só atende a posições $\leq t$.
\end{block}

\vspace{8pt}

\centering
A máscara é o que torna o decoder \textbf{autoregressivo}: cada token é gerado condicionado apenas no que já foi produzido.

\end{frame}

% ============================================================

\begin{frame}[fragile]{Mecânica da Máscara}

\begin{columns}[T]
\column{0.55\textwidth}
\textbf{Como aplicar a máscara:}

\vspace{4pt}

\begin{enumerate}
    \item Calcule os scores $\Qc \Kc^{\top} / \sqrt{d}$ (matriz \texttt{seq $\times$ seq})
    \item Some uma matriz de máscara $M$:
    \begin{itemize}
        \item $M_{ij} = 0$ para $j \leq i$ (válido)
        \item $M_{ij} = -\infty$ para $j > i$ (futuro)
    \end{itemize}
    \item Aplique a softmax: posições com $-\infty$ recebem peso $0$
\end{enumerate}

\vspace{6pt}

\begin{exampleblock}{Resultado}
O modelo nunca olha para tokens à direita: a atenção é estritamente \textbf{causal}.
\end{exampleblock}

\column{0.40\textwidth}
\centering
\begin{tikzpicture}[font=\scriptsize]
\matrix (m) [matrix of nodes,
    nodes={draw, minimum size=0.55cm, anchor=center, inner sep=0pt},
    column sep=-\pgflinewidth, row sep=-\pgflinewidth,
] {
|[fill=lightgreen]| 0 & |[fill=lightred]| {\tiny$-\infty$} & |[fill=lightred]| {\tiny$-\infty$} & |[fill=lightred]| {\tiny$-\infty$} \\
|[fill=lightgreen]| 0 & |[fill=lightgreen]| 0 & |[fill=lightred]| {\tiny$-\infty$} & |[fill=lightred]| {\tiny$-\infty$} \\
|[fill=lightgreen]| 0 & |[fill=lightgreen]| 0 & |[fill=lightgreen]| 0 & |[fill=lightred]| {\tiny$-\infty$} \\
|[fill=lightgreen]| 0 & |[fill=lightgreen]| 0 & |[fill=lightgreen]| 0 & |[fill=lightgreen]| 0 \\
};
\end{tikzpicture}

\vspace{4pt}
\footnotesize Máscara causal triangular.
\end{columns}

\end{frame}

% ============================================================

\archslide{
\textbf{Em foco: cross-attention.}

\vspace{6pt}

O segundo bloco de atenção do decoder consulta a saída do encoder.
}{
\draw[foco] (0.676,0.56) ellipse (0.16 and 0.055);
}

% ============================================================

\begin{frame}{O Decoder: Cross-Attention}

Após o masked self-attention, o decoder aplica um \textbf{segundo} bloco de Multi-Head Attention. Aqui há uma assimetria importante.

\vspace{6pt}

\begin{block}{Cross-Attention (encoder-decoder attention)}
\begin{itemize}
    \item \textcolor{qcolor}{\textbf{Query}}: vem dos estados \textbf{do decoder} (o que estou gerando agora)
    \item \textcolor{kcolor}{\textbf{Key}} e \textcolor{vcolor}{\textbf{Value}}: vêm dos estados \textbf{do encoder} (o que está na entrada)
\end{itemize}
\end{block}

\vspace{6pt}

Em outras palavras: o token alvo ``olha'' para a sequência de origem para decidir o que atender.

\vspace{6pt}

\begin{exampleblock}{Em tradução}
Ao gerar a palavra ``cat'', a query do decoder consulta as keys do encoder e descobre que deve dar mais peso ao value associado ao token original que representa ``gato'' na sequência de origem.
\end{exampleblock}

\end{frame}

% ============================================================

\archslide{
\textbf{Em foco: a cabeça de saída.}

\vspace{6pt}

Uma projeção linear seguida de softmax produz a distribuição sobre o próximo token.
}{
\draw[foco] (0.676,0.86) ellipse (0.155 and 0.065);
}

% ============================================================

\begin{frame}{O Decoder: Feed-Forward e Saída}

\begin{block}{Bloco completo do decoder}
\begin{enumerate}
    \item Masked Multi-Head Attention $\to$ Add \& Norm
    \item Cross-Attention (Q do decoder, K/V do encoder) $\to$ Add \& Norm
    \item Feed Forward $\to$ Add \& Norm
\end{enumerate}
Replicado $N$ vezes (no paper, $N = 6$).
\end{block}

\vspace{6pt}

\begin{exampleblock}{Cabeça final}
Após a pilha de blocos do decoder:
\begin{itemize}
    \item Uma \textbf{Linear} projeta cada vetor de saída para a dimensão do vocabulário $|V|$
    \item Uma \textbf{Softmax} produz uma distribuição de probabilidade sobre o próximo token
\end{itemize}
\end{exampleblock}

\vspace{4pt}

\centering
A geração ocorre token a token: amostramos (ou pegamos o argmax) do próximo token, adicionamos à entrada do decoder, e repetimos.

\end{frame}

% ============================================================
\section{Interpretabilidade}
% ============================================================

\begin{frame}{A Matriz de Atenção como Janela para o Modelo}

Os mecanismos de atenção são \textbf{especialmente fáceis de inspecionar}: podemos plotar os pesos da atenção e ver o que o modelo está ``olhando''.

\vspace{2pt}

\begin{center}
\includegraphics[height=0.50\textheight]{slides/img/04-transformers/voita2019-attention-weights.png}
\end{center}

\vspace{2pt}

\textbf{Exemplo real:} self-attention no encoder de um modelo de tradução. A diagonal predominante mostra que cada token atende sobretudo a si mesmo e aos vizinhos.

\vspace{2pt}
\centering
\footnotesize Figura de Voita et al., 2019. \textit{Analyzing multi-head self-attention.}

\end{frame}

% ============================================================

\begin{frame}{Cabeças Especializadas}

Pesquisas (Voita et al., 2019; entre outros) mostram que algumas \textbf{attention heads} se especializam em fenômenos linguísticos conhecidos.

\vspace{6pt}

\begin{block}{Exemplos de especialização observada}
\begin{itemize}
    \item \textbf{Posicional:} a cabeça atende sempre ao token na mesma posição relativa
    \item \textbf{Concordância verbo-sujeito:} liga formas conjugadas ao seu sujeito
    \item \textbf{Objeto-verbo:} liga objetos diretos ao verbo correspondente
    \item \textbf{Palavras raras:} cabeças que ``buscam'' tokens raros para resolver ambiguidade
\end{itemize}
\end{block}

\vspace{6pt}

\begin{exampleblock}{Conclusão}
Não basta empilhar atenção: as cabeças aprendem \textbf{papéis distintos}. Algumas podem ser podadas com perda mínima, sugerindo redundância controlável.
\end{exampleblock}

\end{frame}

% ============================================================
\section*{Referências e Leituras}
% ============================================================

\begin{frame}{Leituras Recomendadas}
    \leiturasize
    \begin{itemize}
        \item Kazemnejad, A. \textit{Transformer Architecture: The Positional Encoding}. \texttt{\footnotesize kazemnejad.com/blog/transformer\_architecture\_positional\_encoding/}
        \item Alammar, J. \textit{The Illustrated Transformer}. \texttt{\footnotesize jalammar.github.io/illustrated-transformer/}
    \end{itemize}
\end{frame}

\begin{frame}{Referências}

\small

\textbf{Artigos principais:}
\begin{itemize}
    \item Vaswani, A. et al.\ (2017). Attention is All You Need. \textit{Advances in Neural Information Processing Systems}, v.\ 30.
    \item Bahdanau, D., Cho, K., \& Bengio, Y. (2014). Neural Machine Translation by Jointly Learning to Align and Translate. \textit{arXiv:1409.0473}.
    \item Luong, M.-T., Pham, H., \& Manning, C. D. (2015). Effective Approaches to Attention-based Neural Machine Translation. \textit{arXiv:1508.04025}.
\end{itemize}

\vspace{4pt}

\textbf{Análise e interpretabilidade:}
\begin{itemize}
    \item Voita, E., Talbot, D., Moiseev, F., Sennrich, R., \& Titov, I. (2019). Analyzing Multi-Head Self-Attention: Specialized Heads Do the Heavy Lifting, the Rest Can Be Pruned. \textit{arXiv:1905.09418}.
\end{itemize}

\end{frame}

\end{document}
